\begin{abstract}[Resumen]
Este Trabajo de Fin de Máster presenta el diseño, implementación y evaluación de un asistente institucional basado en \textit{Retrieval-Augmented Generation} (RAG) e integración de herramientas. La arquitectura propuesta combina recuperación documental sobre un corpus de fuentes institucionales, mejora del \textit{ranking} de evidencias mediante recuperación híbrida y \textit{reranking}, orquestación del flujo de interacción basado en agente y generación y envío supervisado de correos electrónicos.

El sistema se evalúa mediante un \textit{ablation study} con cuatro variantes progresivas sobre un dataset de 71 preguntas anotadas manualmente sobre documentación pública de UNIR, analizando el impacto de cada componente sobre métricas de recuperación (Hit@1, Hit@3, MRR) y de generación (Fact-Coverage, ROUGE-L, BERTScore).

Los resultados muestran que la recuperación híbrida y el \textit{reranking} incrementan la cobertura factual un 35,5\,\% respecto al \textit{baseline} (de 0.389 a 0.527) y mejoran la cobertura en el top-3 (Hit@3 = 0.918). El orquestador aporta valor en términos de control del flujo, abstención ante preguntas no respondibles y activación de herramientas accionables, con una calidad semántica equivalente a la variante determinista (BERTScore 0.648 frente a 0.656). En conjunto, los resultados sugieren la utilidad de una arquitectura modular, trazable y supervisada para asistentes institucionales privados.

\textbf{Palabras clave:} Modelos de lenguaje de gran tamaño, Retrieval-Augmented Generation, recuperación híbrida, reranking, orquestación, uso de herramientas, asistente institucional.
\end{abstract}


\pagebreak
%Ingles
\begin{abstract}[Abstract]
This Master's Thesis presents the design, implementation and evaluation of an institutional assistant based on Retrieval-Augmented Generation (RAG) and tool integration. The proposed architecture combines document retrieval over an institutional corpus, evidence ranking improvement through hybrid retrieval and reranking, agent-based interaction flow orchestration, and supervised generation and sending of emails.

The system is evaluated through a four-variant ablation study on a manually annotated dataset of 71 questions over public UNIR documentation, measuring the impact of each component on retrieval metrics (Hit@1, Hit@3, MRR) and generation metrics (Fact-Coverage, ROUGE-L, BERTScore).

Results show that hybrid retrieval and reranking improve fact coverage by 35.5\% over the baseline (from 0.389 to 0.527) and increase top-3 coverage (Hit@3 = 0.918). The orchestrator contributes flow control, abstention on unanswerable questions and activation of actionable tools, with semantic quality equivalent to the deterministic variant (BERTScore 0.648 vs 0.656). Overall, the results suggest the usefulness of a modular, traceable and supervised architecture for private institutional assistants.

\textbf{Keywords:} Large Language Models, Retrieval-Augmented Generation, hybrid retrieval, reranking, orchestration, tool use, institutional assistant.
\end{abstract}